\section*{第 \ref{ch:g2:teeth-and-eating-well} 章——牙齿与健康饮食}

\begin{solution}{exo:g2:teeth-and-eating-well:1}
\omterm{def:g2:teeth-and-eating-well:kinds}{门齿}用来切；\omterm{def:g2:teeth-and-eating-well:kinds}{犬齿}用来撕；\omterm{def:g2:teeth-and-eating-well:kinds}{臼齿}用来磨。
\end{solution}

\begin{solution}{exo:g2:teeth-and-eating-well:2}
\omterm{def:g2:teeth-and-eating-well:kinds}{门齿}位于前面（\omterm{def:g2:teeth-and-eating-well:kinds}{犬齿}在其两侧的角上）；
\omterm{def:g2:teeth-and-eating-well:kinds}{臼齿}位于后面。
\end{solution}

\begin{solution}{exo:g2:teeth-and-eating-well:3}
\omterm{def:g2:teeth-and-eating-well:kinds}{门齿}切下一口食物；\omterm{def:g2:teeth-and-eating-well:kinds}{犬齿}帮助抓住并把食物拉下来；
舌头把食物送回；\omterm{def:g2:teeth-and-eating-well:kinds}{臼齿}把它磨成糊状；
然后被咽下。
\end{solution}

\begin{solution}{exo:g2:teeth-and-eating-well:4}
\omterm{prop:g2:teeth-and-eating-well:sets}{乳牙}共有 $20$ 颗。大约从六岁起，它们开始松动，
一颗接一颗地脱落，由下方生长中的
\omterm{prop:g2:teeth-and-eating-well:sets}{恒牙}顶出。
\end{solution}

\begin{solution}{exo:g2:teeth-and-eating-well:5}
因为\omterm{prop:g2:teeth-and-eating-well:sets}{恒牙}是最后一套：脱落的\omterm{prop:g2:teeth-and-eating-well:sets}{乳牙}会被\omterm{prop:g2:teeth-and-eating-well:sets}{恒牙}替换，
但脱落的\omterm{prop:g2:teeth-and-eating-well:sets}{恒牙}永远不会再长出。整套\omterm{prop:g2:teeth-and-eating-well:sets}{恒牙}
必须用一辈子。
\end{solution}

\begin{solution}{exo:g2:teeth-and-eating-well:6}
水果和蔬菜；面包和谷物；奶制品（牛奶、酸奶、
奶酪）；肉、鱼或蛋；零食只宜少量、
偶尔吃。身体真正需要的唯一饮品是水。
\end{solution}

\begin{solution}{exo:g2:teeth-and-eating-well:7}
苹果提供维生素，面包提供持久的能量，水则必不可少
——这才是真正的营养，也让牙齿得到实实在在的咀嚼锻炼。
零食 B 主要是糖，会喂养侵蚀牙齿的\omterm{def:g1:healthy-habits:germs}{病菌}，
而非有益的食物。
\end{solution}

\begin{solution}{exo:g2:teeth-and-eating-well:8}
开放性答案——对着镜子观察：前面是扁平、边缘锋利的\omterm{def:g2:teeth-and-eating-well:kinds}{门齿}，
每个角上是一颗尖尖的\omterm{def:g2:teeth-and-eating-well:kinds}{犬齿}，后面是宽大、凹凸不平的\omterm{def:g2:teeth-and-eating-well:kinds}{臼齿}，
与图中一致。
\end{solution}

\begin{solution}{exo:g2:teeth-and-eating-well:9}
扁平的切割用\omterm{def:g2:teeth-and-eating-well:kinds}{门齿}和磨碎用\omterm{def:g2:teeth-and-eating-well:kinds}{臼齿}是\omterm{def:g2:what-animals-eat:kinds}{植食性动物}的工具式样；
尖尖的\omterm{def:g2:teeth-and-eating-well:kinds}{犬齿}是\omterm{def:g2:what-animals-eat:kinds}{肉食性动物}的式样。两者兼有，意味着既能
吃\omterm{def:g1:plants-around-us:plant}{植物}也能吃肉的“菜单”：是\omterm{def:g2:what-animals-eat:kinds}{杂食性动物}的口腔，
正如该章中黑鸟和野猪的食谱那样。
\end{solution}

\begin{solution}{exo:g2:teeth-and-eating-well:10}
缺少的有：水果和蔬菜（一整天一点都没有！）以及作为日常饮品的水；
肉、鱼或蛋也没有。吃得太多的是：
含糖食品——糖果、汽水、饼干。改进办法：早餐和加餐吃水果，
午餐或晚餐吃蔬菜和一些鱼或蛋，
用水代替汽水，零食限制在一种。
\end{solution}
